\lecture{4}{Wed 07 Sep 2022 11:38}{Finite, Countable and Uncountable Sets}

\subsection{Finite and countable sets}
\begin{definition}
	We say that two sets $A,B$ are \textit{equivalent} (or have the same number of elements) if there is a bijection \[
	\varphi : A\to B
.\] We write $A\sim B$ for this equivalence relation.

\end{definition}

\begin{proposition}
	Reflexivity: $A\sim A$ ($\iota_A)$

	Symmetry: If $A\sim B$ then $B\sim A$. ($\varphi^{-1}:B\to A$ is a bijection)

	Transitivirty: If $A\sim B$ and $B\sim C$ then $A\sim C$. We compose the bijections for the proof.
\end{proposition}


Define $S_n = \{1,2,3,\ldots,n\} $ be the initial segment of $\mathbb{N}$, and let $S_0 = \emptyset $. We say that a set $E$ is finite if there is a bijection $  S_n\to E$ for some $n \in \mathbb{N}$. In other words, $E \sim S_n$.

We say that $E$ is \textit{denumerable} if there is abijection $\mathbb{N}\to E$, i.e., $E\sim \mathbb{N}$. If the set is not finite, it is called infinite. 


Note $\mathbb{N}$ is infinite: does not exist bijection $S_n \to \mathbb{N}$.

But there is a bijection between $\mathbb{N}$ and its proper subset.

We say set is \textit{countable} if it is finite or denumerable. The infinite sets that are not countable are called \textit{uncountable}.

\begin{proposition}
	$S_{n}\sim S_m \iff n = m$ for $n, m \in \{0\} \cup \mathbb{N}\cup \{\infty \} $.
\end{proposition}
\begin{proof}
	(partial) Show if $n \in \mathbb{N}$ then $S_n$ is not equivalent to $\mathbb{N}$. Suppose the countary and \[
		S_n \to \mathbb{N}
	\]
	is a bijection. Then look at $N = \max \{\varphi(1), \ldots, \varphi(n)\} + 1$ . Now  $N > \varphi(k)$ for any $k$. Therefore $N \not\in \varphi(S_n)$. Therefore $\varphi(S_n) \neq  \mathbb{N}$.
\end{proof}

\begin{proposition}
	\begin{enumerate}
		\item Any subset of a finite set is finite.
		\item Any infinite subset of a denumerable set is denumerable.
		\item Any subset of a countable set is countable.
	\end{enumerate}
	Note that (1) and (2) imply (3).
\end{proposition}
\begin{proof}
	(of (2)). Enough to show that if $E \subset \mathbb{N}$ then $E$ is countable. We use the Well-ordering property of $\mathbb{N}$: any non-empty subset has a minimum element.

	$E$ is infinite, so $E$ is nonempty. Therefore exists minimal $n_1 \in E$. $E$ is not finite, so $E \setminus \{n_1\} $ is nonempty. Then there exists $n_2 \in E \setminus \{n_1\} $.

	Continue in this fashion. If we constructed $n_1, \ldots, n_k \in E$, such that $n_1 < n_2 <\ldots<n_k$, then $E\setminus \{n_1, \ldots, n_k\} $ is nonempty since $E$ is infinite. Therefore there exists $n_{k+1}$ be the smallest element of $E\setminus \{n_1, \ldots, n_k\} $.

\begin{claim}
	$E = \{n_1, \ldots, n_k, \ldots\} $
\end{claim}
\begin{proof}


\begin{lemma}
	$n_k \ge k$ for any $k \in \mathbb{N}$.
	proof by induction.
\end{lemma}

Note that by construction, $n_1, \ldots, n_k, \ldots \subset E$. We need to show that other direction.

Let $m \in E$. By Archimedean property, $\exists k \in \mathbb{N}$ s.t. $m<k\le n_k$.

Let $k_0$ be the smallest number such that  $m < n_{k_0}$. This means that $n_{k_0-1}\le m$. If $n_{k_0-1}\le m$ then we are done. If $n_{k_0-1}< m$, then $m \in E\setminus \{n_1, .., n_{k_{0}-1}\} $, hence $m < n_{k_0}$, a contradiction with construction of $n_{k_0}$.
\end{proof}
\end{proof}

\begin{proposition}
	\begin{enumerate}
		\item Finite union of finite sets is finite.
		\item Denumerable union of denumerable sets is denumerable.
		\item Countable union of countable sets is countable.
	\end{enumerate}
\end{proposition}
To prove (2) of this proposition, we need
\begin{proof}
	
\begin{lemma}
	$\mathbb{N}\times \mathbb{N}\sim \mathbb{N}$.
\end{lemma}
\begin{proof}
	Use diagonal construction.
\end{proof}


Let
\[
	E = \bigcup_{n \in \mathbb{N}} E_n
\]
Note that $E_n $ might overlap. We define
\[
	\varphi_n: \mathbb{N} \to  E_n
.\] 
and define 

\begin{align*}
	\varphi: \mathbb{N}\times \mathbb{N} &\longrightarrow \bigcup_{n \in \mathbb{N}} E_n = E \\
	n,m &\longmapsto \varphi(n,m) =  \varphi_n(m)
.\end{align*}

We may view $\varphi$ as $\varphi: \mathbb{N} \to E$ due to isomorphism between $\mathbb{N}\times \mathbb{N}$ and $\mathbb{N}$.

Now for $e \in E$, pick the smallest element in $\varphi^{-1}(e)$, call it $f(e)$. $f: E \to  \mathbb{N}$ and $\varphi \circ f = 1_E$. Hence $E \sim f(E) \subset \mathbb{N}$, hence $E$ is countable and infinite, hence $E$ is denumerable.

\end{proof}

\begin{corollary}
	$\mathbb{Q}$ is countable, since \[
		\mathbb{Q} = \bigcup_{n \in \mathbb{N}} \frac{1}{n} \mathbb{Z}
	.\] 
\end{corollary}

\begin{theorem}[Cantor]
	$\mathcal{P}$ is uncountable, hence $F$ is uncountable.
\end{theorem}
\begin{proof}
	By Cantor's diagonal argument.
	Suppose $\mathcal{P}$ is countable, then \[
		\mathcal{P} = \{\alpha^{(1)}, \alpha^{(2)}, \ldots,\alpha^{(n)},\ldots \} 
	.\] 
	and \[
		\alpha_1^{(1)} = \{\alpha_1^{(1)}, \alpha_2^{(1)}, \ldots,\alpha_1^{(n)},\ldots \} 
	.\] 

	We construct $\beta \in \mathcal{P}$ such that $\beta \neq  \alpha^{(n)}$ for any $n$.
	Take $\beta_1 = \overline \alpha_1^{(1)} = 1- \alpha_1^{(1)}$, and $\beta_n = \overline \alpha_n^{(n)} = 1- \alpha_n^{(n)}$.

	Claim: $\beta$ does not equal to $\alpha^{(n)}$ for any $n$. Proof: Note that $\beta$ and $\alpha^{(n)}$ differ at the $n$-th element. 

	Now we have a contradiction since $\beta \in \mathcal{P}$. Hence $\mathcal{P}$ cannot be countable.
\end{proof}

\begin{corollary}
	$\mathbb{R}$ is uncountable. $\mathbb{R}\setminus \mathbb{Q}$ is uncountable.
\end{corollary}
